%% kapa: exam=1 round=30 year=2019 period=1 form=A subject=real_estate subject_ko=부동산학원론 seq=24 number=104
%% source: Q-Net 감정평가사 기출문제 (cst00309) artlSeq=5205205 file=extracted/1008/5205205/2019년 제30회 감정평가사 1차시험 문제지/2019년 제30회 감정평가사 1차시험 문제지/2019년도 제30회 감정평가사 1차 1교시 A형.hwp
%% conversion: kordoc 4.24.51 -> tools/gen1.py (auto); review=proofread-sonnet; flags=none
\begin{question}{104}
감정평가에 관한 규칙상 시산가액 조정에 관한 설명으로 옳지 않은 것은?
\begin{choices}
  \item 평가대상물건별로 정한 감정평가방법을 적용하여 산정한 가액을 시산가액이라 한다.
  \item 평가대상물건의 시산가액은 감정평가 3방식 중 다른 감정평가방식에 속하는 하나 이상의 감정평가방법으로 산정한 시산가액과 비교하여 합리성을 검토하여야 한다.
  \item 시산가액 조정시 공시지가기준법과 거래사례비교법은 같은 감정평가방식으로 본다.
  \item 대상물건의 특성 등으로 인하여 다른 감정평가방법을 적용하는 것이 곤란하거나 불필요한 경우에는 시산가액 조정을 생략할 수 있다.
  \item 산출한 시산가액의 합리성이 없다고 판단되는 경우에는 주된 방법 및 다른 감정평가방법으로 산출한 시산가액을 조정하여 감정평가액을 결정할 수 있다.
\end{choices}
\end{question}
\answer{3}
